\documentclass{article}
\usepackage[margin=1in]{geometry}
\usepackage{hyperref}
\usepackage{parskip}
\usepackage{lmodern}

\pagestyle{empty}

\begin{document}

\begin{flushleft}
    {\Large Alexandre Ait Ettajer} \\
    \href{mailto:alexandre.aitettajer@gmail.com}{alexandre.aitettajer@gmail.com} \quad +1(281) 435-9481 \quad Boston, MA \\
    GitHub: \href{https://github.com/aiteta}{Aiteta} \quad LinkedIn: \href{https://www.linkedin.com/in/alexandre-ait-ettajer-65740b170/}{alexandre-ait-ettajer}
\end{flushleft}

\vspace{1em}

Dear Multigres team,

I am applying for the Multigres Engineer role. I am a software engineer with four years at MathWorks and an M.S. in Computer Science (distributed and parallel systems, UT Austin), and I build distributed systems in Go — including a fault-tolerant distributed file system that uses Paxos consensus for replication across nodes. Consensus protocols are not a keyword on my resume; they are the core mechanism of the hardest system I have built.

At MathWorks, I steered a working group on distributed visualization architectures across teams and rebuilt the client-server messaging architecture behind MATLAB visualizations so downstream teams could ship their own cloud-native visualizations. I also designed tabular data support for core graphics, integrating SQL-style datasets — experience that maps directly onto working with Postgres at the storage layer.

What draws me to Supabase is that Multigres sits exactly at the intersection I want to work in: sharded, replicated Postgres with real consensus underneath, built in the open, for a global remote team. As an EU citizen fluent in French and English, I am set up well for distributed async collaboration across time zones.

I would welcome the chance to talk about how I think about replication, fault tolerance, and building databases people rely on.

Sincerely, \\
Alexandre Ait Ettajer

\end{document}
